\documentclass[11pt]{article}
\usepackage{amsmath, amssymb, amsfonts, geometry}
\geometry{margin=1in}

\title{Supplementary Information: Information Distortion, Effective Sample Size, and Corrected Covariance in Sequential Biophysical Measurements}
\author{}
\date{}

\begin{document}

\maketitle

\section*{1. Log-Likelihood and Score Definitions}
Consider a trajectory of $T$ sequential measurement intervals (e.g., time steps in a patch-clamp recording or frames in a single-molecule fluorescence trace). [cite_start]The total log-likelihood is defined as[cite: 2, 3]:
\begin{equation}
\ell(\theta) = \sum_{t=1}^T \ell_t(\theta),
\end{equation}
[cite_start]with the total score $S(\theta)$, representing the gradient of the log-likelihood, defined as the sum of per-interval scores $s_t(\theta)$[cite: 3]:
\begin{equation}
S(\theta) = \sum_{t=1}^T s_t(\theta), \qquad s_t(\theta) = \nabla_\theta \ell_t(\theta).
\end{equation}

\section*{2. Model Curvature: Gaussian Fisher Information (GFI)}
[cite_start]The model approximates the predictive likelihood at each interval as Gaussian: $y_t \sim \mathcal{N}(\mu_t(\theta), \sigma_t^2(\theta))$[cite: 5]. [cite_start]The total Gaussian Fisher Information $\mathbf{H}(\theta)$ represents the model's predicted sensitivity or curvature[cite: 7]:
\begin{equation}
\mathbf{H}(\theta) = \sum_{t=1}^T \mathbb{E}\left[\mathbf{I}_t^{\mathrm{G}}(\theta)\right] = - \mathbb{E}\left[\nabla_\theta^2 \ell(\theta)\right].
\end{equation}
In a biophysical context, $\mathbf{H}$ represents the expected precision of parameter estimates based on the assumed physical model (e.g., the sharpness of an energy landscape).

\section*{3. Measuring Fluctuations: Sample vs. Total Covariance}
To diagnose model failure, we compare the predicted curvature $\mathbf{H}$ against the actual observed fluctuations of the score. [cite_start]We distinguish between local noise and trajectory-wide dependencies[cite: 29, 31]:

\begin{itemize}
    [cite_start]\item \textbf{Sample Score Covariance ($\mathbf{J}_{\mathrm{sample}}$):} The sum of the variances of individual intervals[cite: 35]. This captures the ``local'' volatility of the data, as if each measurement were an independent sample:
    \begin{equation}
    \mathbf{J}_{\mathrm{sample}} = \sum_{t=1}^T \mathrm{Cov}(s_t).
    \end{equation}
    [cite_start]\item \textbf{Total Score Covariance ($\mathbf{J}_{\mathrm{total}}$):} The covariance of the total score sum, which captures the ``global'' volatility[cite: 8]. It includes both the local noise and all cross-interval temporal correlations:
    \begin{equation}
    \mathbf{J}_{\mathrm{total}} = \mathrm{Cov}(S(\theta)) = \mathbf{J}_{\mathrm{sample}} + 2\sum_{i<j}\mathrm{Cov}(s_i,s_j).
    \end{equation}
\end{itemize}

\section*{4. Derivation of the Distortion-Corrected Covariance (DCC)}
The standard error of parameter estimates $\hat{\theta}$ is often underestimated in biophysics due to unmodeled correlations. We derive a robust correction by performing a first-order Taylor expansion of the score function around the true parameters $\theta_0$:
\begin{equation}
S(\hat{\theta}) \approx S(\theta_0) + \nabla_\theta S(\theta_0) (\hat{\theta} - \theta_0).
\end{equation}
At the maximum likelihood estimate, the score is zero ($S(\hat{\theta}) = 0$). Replacing the empirical Hessian with its expectation ($-\mathbf{H}$), we solve for the estimation error:
\begin{equation}
(\hat{\theta} - \theta_0) \approx \mathbf{H}^{-1} S(\theta_0).
\end{equation}
The true covariance of the parameter estimates, $\Sigma_{\mathrm{DCC}}$, is the expectation of the squared error:
\begin{equation}
\Sigma_{\mathrm{DCC}} = \mathbb{E}\left[(\hat{\theta} - \theta_0)(\hat{\theta} - \theta_0)^\top\right] \approx \mathbf{H}^{-1} \mathbb{E}[S S^\top] \mathbf{H}^{-1}.
\end{equation}
[cite_start]Since $\mathbf{J}_{\mathrm{total}} = \mathbb{E}[S S^\top]$ (assuming a centered score), we arrive at the corrected covariance[cite: 11, 27]:
\begin{equation}
\Sigma_{\mathrm{DCC}} = \mathbf{H}^{-1} \mathbf{J}_{\mathrm{total}} \mathbf{H}^{-1}.
\end{equation}
This ``sandwich'' structure ensures that if the actual score fluctuations ($\mathbf{J}_{\mathrm{total}}$) are larger than the model-predicted curvature ($\mathbf{H}$), the parameter uncertainty is appropriately inflated.

\section*{5. Decomposition of Information Distortion}
[cite_start]The total Information Distortion Matrix ($\mathbf{C}$) quantifies the mismatch between the model's assumed geometry and the data's actual behavior[cite: 9, 28]:
\begin{equation}
\mathbf{C} = \mathbf{H}^{-1/2} \mathbf{J}_{\mathrm{total}} \mathbf{H}^{-1/2}.
\end{equation}
[cite_start]We decompose this distortion into two physically distinct components[cite: 19, 42]:

\subsection*{5.1 Sample Distortion ($\mathbf{C}_{\mathrm{sample}}$)}
[cite_start]This measures the instantaneous mismatch between the observed volatility of a single measurement and the model's predicted curvature[cite: 40, 41]:
\begin{equation}
\mathbf{C}_{\mathrm{sample}} = \mathbf{H}^{-1/2} \mathbf{J}_{\mathrm{sample}} \mathbf{H}^{-1/2}.
\end{equation}
If $\mathbf{C}_{\mathrm{sample}} \neq \mathbf{I}$, the model's physical assumptions (e.g., noise levels or state-space topology) are incorrect even at the single-sample level.

\subsection*{5.2 Correlation Distortion ($\mathbf{R}$)}
[cite_start]This measures the inflation of uncertainty due to temporal memory effects (e.g., slow conformational changes or diffusion)[cite: 35, 36]:
\begin{equation}
\mathbf{R} = \mathbf{J}_{\mathrm{sample}}^{-1/2} \mathbf{J}_{\mathrm{total}} \mathbf{J}_{\mathrm{sample}}^{-1/2}.
\end{equation}
If the data is temporally uncorrelated, $\mathbf{R} = \mathbf{I}$. [cite_start]The full distortion is thus a product of modeling errors and temporal dependencies[cite: 21, 42]:
\begin{equation}
\mathbf{C} = \mathbf{C}_{\mathrm{sample}}^{1/2} \mathbf{R} \mathbf{C}_{\mathrm{sample}}^{1/2}.
\end{equation}

\section*{6. Bias and Effective Sample Size}
[cite_start]Together, these quantities distinguish between different failure modes of a biophysical model[cite: 28, 47]:
\begin{itemize}
    [cite_start]\item \textbf{Effective Sample Size:} If distortion is purely due to isotropic temporal correlation, then $\mathbf{C} \approx \kappa \mathbf{I}$, and the effective sample size is $T_{\mathrm{eff}} = T/\kappa$[cite: 15, 34].
    [cite_start]\item \textbf{Distortion-Induced Bias (DIB):} If the expected score is non-zero ($\mathbb{E}[S] = \mathbf{g}$), it indicates a systematic drift or bias in the model[cite: 12, 46]:
    \begin{equation}
    \mathbf{b}_{\mathrm{DIB}} = \mathbf{H}^{-1} \mathbf{g}.
    \end{equation}
\end{itemize}

\end{document}